\cotefile{161-3}
\scanpage{4}
% L1
\uline{Produits de topos}

% L2–3
$\Pi(E,F)$ \unsure{$=$|$\simeq$} faisceaux sur $E$ à valeurs dans $F$ i.e. foncteurs $E^o\to F$ \deleted{\unsure{respectant}}commutant aux $\varprojlim$, i.e. qui ont un adjoint à gauche

% L4–5
$\simeq$ faisceaux sur $F$ à valeurs dans $E$, i.e foncteurs $F^o\unsure{\longrightarrow|-}E$ commutant aux $\varprojlim$, i.e. qui ont un adjoint à g.

% L6–7
$\simeq$ foncteurs $E^o\times F^o\longrightarrow\mathrm{Ens}$ qui commutent \uline{en chaque} \uline{variable} aux $\varprojlim$ quelc.

% L8-left
\margin{Si $E\cong\widetilde C$}

% L8–13
$\simeq$ faisceaux sur $C$ à valeurs dans $F$, i.e. foncteurs $\varphi\colon C^o\longrightarrow F$ tels que pour tt $R\hookrightarrow X$ couvrant dans $C$, on a $\varphi(x)\xrightarrow{\sim}\varprojlim_{\unsure{C/R|c/R}}\varphi(z)$, i.e. tels que $\forall\unsure{y|Y}\in F$, $x\mapsto\mathrm{Hom}(\varphi(x),y)$ sur $C$ est un faisceau, \inserted{(\unsure{ou autres|ou plutôt} \unsure{repr.|resp.|syst.} de flèches)} i.e. foncteurs $\varphi^o\colon C\to F^o$ qui sont ``\uline{continus}'' [i.e. se prolongent en foncteurs $\unsure{\overline C|\widetilde C}\longrightarrow F^o$ commutant aux $\varinjlim$ quelcon.]

% L14-left
\margin{Si $F\cong\widetilde D$}

% L14–16
$\simeq$ faisceaux sur $D$ à valeurs dans $E$, i.e. foncteurs $\psi\colon D^o\longrightarrow E$ tels que --, i.e. foncteurs ``continus'' $\psi^o\colon D\to E^o$

% L17–19
$\simeq$ foncteurs $C^o\times D^o\longrightarrow(\mathrm{Ens})$ qui sont des ``bifaisceaux'' i.e. tels que ce soient des faisceaux séparément en chaque variable

% L20
\deleted{\unsure{Soient} $u\colon E\longrightarrow E'$, $v\colon F\longrightarrow F'$ \unsure{d}}

% L21
\uline{Prop. \textcircled{1}} \quad $\Pi(E,F)$ est un topos

% L22
On peut supposer $E=\widetilde C$, $F=\widetilde D$, $C$ et $D$ petits sites stables par $\varprojlim$ finies.

% L23
\unsure{a)} Cas des topologies discrètes : $\pi(\widehat C,\widehat D)\cong\widehat{C\times D}$ : c'est un topos

% L24–25
b) Plus généralement, $\pi(\widehat C,\widetilde D)\simeq\underline{\mathrm{Hom}}\,(C^o,\widetilde D)$ est un topos, \deleted{\unsure{et} de \unsure{même}} $\pi(\widetilde C,\widehat D)\simeq\underline{\mathrm{Hom}}\,(D^o,\widetilde C)$.

% L26–29
\deleted{\unsure{Le|Et}} L'inclusion naturelle $\pi(\widehat C,\widetilde D)\hookrightarrow\Pi(\widehat C,\widehat D)$ \inserted{($\varphi\mapsto i_D\circ\varphi$, où $i_D\colon\widetilde D\longrightarrow\widehat D$ l'inclusion)} est le foncteur image directe pour un morphisme de topos, dont le foncteur image inverse \deleted{\unsure{est}} $\pi(\widehat C,\widehat D)\longrightarrow\pi(\widehat C,\widetilde D)$ est donné par $\varphi\longmapsto a_D\circ\varphi$, où $a_D\colon\widehat D\longrightarrow\widetilde D$ est l'adjoint à g. de $i_D$.

% L30–32
Donc $\pi(\widehat C,\widetilde D)$ est un sous-topos (str. plein) dans $\pi(\unsure{\widehat C|C},\widehat D)$ et pour la raison symétrique $\pi(\widetilde C,\widehat D)$ aussi. Mais $\pi(\widetilde C,\widetilde D)$ est l'intersection de ces deux, donc c'est un \inserted{sous-}topos de $\pi(\widehat C,\widehat D)$

% foot
% archive sheet number: 3
\scanpage{5}
% L1
a fortiori un topos.

% L2
\deleted{\illegible{1}}

% L3
On va définir, pour deux morphismes de topos

% L4
\[
\deleted{\unsure{f}}u\colon E\longrightarrow E',\qquad \deleted{\unsure{g}}v\colon F\to F'
\]

% L5
un morphisme de topos

% L6
\[
\Pi(u,v)\colon\pi(E,F)\longrightarrow\pi(E',F')
\]

% L7–8
dont l'interprétation intrinsèque symétrique est ainsi, en termes des ``bifaisceaux'' $E^o\times F^o\longrightarrow(\mathrm{Ens})$, et $E'{}^o\times F'{}^o\longrightarrow(\mathrm{Ens})$,

% L9–11
\[
\begin{aligned}
&\deleted{\unsure{\pi(u,v)^*(\varphi')(x,y)=\varphi'(u_*(x),v_*(y))}}\\
&\Pi(u,v)_{\ast}^{\deleted{\ast}}(\varphi)(X',Y')=\varphi(u^*X',v^*Y')\quad:\pi(E,F)\longrightarrow\pi(E',F')\\
&\pi(u,v)^*(\varphi)(X,Y)\overset{?}{=}
\end{aligned}
\]

% L12–13
Pour prouver que \deleted{\unsure{c'est associé}} foncteur précédent est associé à un morphisme de topos, nous sommes ramenés, en \deleted{\unsure{fa|le}} factorisant en

% L14
\[
\pi(E,F)\longrightarrow\pi(E,F')\longrightarrow\pi(E',F')
\]

% L15–16
au cas où l'un des deux morphismes $u,v$ est l'identité ; soit $u$, donc à étudier

% L17
\[
\pi(\mathrm{id}_E,v)_*\colon\pi(E,F)\longrightarrow\pi(E,F')
\]

% L18–21
On peut supposer \deleted{$v\colon F\to F'$ défini par un morphisme de \inserted{petits} sites $v_0^*\colon C'\to C$, $C,C'$ stables par \deleted{\unsure{produits}} $\varprojlim$ finies, leur top. moins fine que la top can. \deleted{Donc} On a \deleted{\unsure{évidemment}} alors un diagramme commutatif \unsure{$\pi(\mathrm{id}_E,v)_*$} de topos}

% L22–25-left
\diagram{\begin{tikzcd}[ampersand replacement=\&,column sep=large,execute at end picture={\draw (current bounding box.north west) -- (current bounding box.south);}]
\deleted{F}\arrow[r,"{\deleted{\unsure{v_0|\widetilde v_0}}}"] \& \deleted{F'}\\
\deleted{\widetilde C=F}\arrow[r,"{\deleted{\widetilde v_0}}"]\arrow[d,hook] \& \deleted{\widetilde C'=F'}\arrow[d,hook]\\
\deleted{\widehat C}\arrow[r,"{\deleted{\widehat v_0}}"] \& \deleted{\unsure{\widehat C'|C'}}
\end{tikzcd}}

% L26-left
\deleted{et par suite un diagr\unsure{.}}

% L22–25-right
$E=\widetilde C$, et \deleted{\unsure{on a}} alors $\pi(\mathrm{id}_E,v)_*$ s'insère dans un diagramme \inserted{\unsure{ess.}} commutatif

% L26–29
\diagram{\begin{tikzcd}[ampersand replacement=\&,column sep=huge,row sep=large]
\pi(\widetilde C,F)\arrow[r,"{\pi(\mathrm{id}_{\widetilde C},v)_*}"]\arrow[d,hook,"{\pi(i_C,\mathrm{id}_F)_*}"'] \& \pi(\widetilde C,F')\arrow[d,hook,"{\pi(i_C,\mathrm{id}_{F'})_*}"]\\
\pi(\widehat C,F)\arrow[r,"{\pi(\mathrm{id}_{\widehat C},v)_*}"'] \& \pi(\widehat C,F')
\end{tikzcd}}

% foot
% archive sheet number: 4
\scanpage{6}
% L1–5
Nous avons déjà \deleted{\unsure{sur}} \inserted{dans la dém de} prop. 1 vérifié que les flèches verticales sont des inclusions de sous-topos [inclusion d'un sous-topos \unsure{ensembliste|ensemblistique|intersection} dans un des facteurs] donc pour prouver que $\Pi(\mathrm{id}_{\widetilde C},v)_*$ est le \inserted{foncteur} image directe pour un morphisme de topos, il suffit de vérifier qu'il en est ainsi de $\Pi(\mathrm{id}_{\widehat C},F')_*$. En effet,

% L6
On aura alors un morphisme de topos

% L7
\diagram{\begin{tikzcd}[ampersand replacement=\&,row sep=small,column sep=small]
\pi(\widetilde C,F)\arrow[rr]\arrow[dr] \& \& \pi(\widehat C,F')\\
\& \pi(\widehat C,F)\arrow[ur] \&
\end{tikzcd}}

% L7–11
, et on voit (en vérifiant sur les foncteurs images directs) qu'il se factorise par le sous topos $\pi(\unsure{\overline C|\widetilde C},F')$, le foncteur image directe associé au \unsure{morphisme} \illegible{1} topos obtenu $\pi(\widetilde C,F)\longrightarrow\pi(\widetilde C,F')$ étant $\pi(\mathrm{id}_{\widetilde C},v)_*$.

% L12
Or $\pi(\mathrm{id}_{\widehat C},\unsure{v)_*}$ \unsure{peut} s'identifier au foncteur

% L13
\[
\underline{\mathrm{Hom}}\,(C,F)\unsure{\longrightarrow|-}\underline{\mathrm{Hom}}\,(C,F')
\]

% L14–19
défini par \deleted{\unsure{$\varphi\mapsto\varphi\circ v^*$}, dont l'adjoint est $\varphi'\mapsto\unsure{\varphi_*}v_*\circ\varphi'$,} $\varphi\longmapsto v_*\circ\varphi$ [car le bifoncteur associé à $\varphi$ est \deleted{\unsure{$\varphi$}} $(X,Y)\to\mathrm{Hom}_F(Y,\varphi(X))$, qu'on \unsure{remplace|compose|transforme} par \deleted{\unsure{transport|Hom}} $(X,Y')\unsure{\longrightarrow|-}\mathrm{Hom}_F(v^*Y',\varphi(X))\simeq\mathrm{Hom}_{F'}(Y',(v_*\varphi)(X))$, et posant $\varphi'=v_*\varphi\colon C\longrightarrow F'$, c'est bien le bifoncteur sur $C\times F'$ associé à $\varphi'$], dont l'adjoint à droite est $\varphi\longmapsto v^*\circ\varphi$, or celui-ci est bien exact à g. \qquad cqfd.

% L20–21
Prenant $v\colon F\longrightarrow P$ (topos ponctuel $\simeq\widehat e$, $e$ catégorie ponctuelle) et notant que $\pi(E,P)\cong E$, on trouve un morphisme can.

% L22
\[
\pi(E,F)\longrightarrow E
\]

% L23
et on trouve de $=$ un morphisme can. de topos

% L24
\[
\pi(E,F)\longrightarrow F
\]

% L25
d'où un morphisme de topos

% L26
\[
\pi(E,F)\longrightarrow E\mathbin{\underset{\mathrm{top}}{\overset{2}{\times}}}F
\]

% L27–28
évidemment 2-fonctoriel en $E,F$. On veut prouver que c'est une \deleted{\unsure{trans}} équivalence, i.e. que $\pi(E,F)$ ``est'' un 2-produit des

% foot
% archive sheet number: 5
\scanpage{7}
% L1–2
topos $E,F$. On suppose comme dans la dém. de prop 1 $E=\widetilde C$, $F=\widetilde D$, \deleted{\unsure{donc}} on a vu que l'on a alors

% L3–4
\[
\begin{aligned}
\pi(E,F)=\pi(\unsure{\overline C|\widetilde C},\widetilde D)
 &=\deleted{\unsure{\pi}}\pi(\widetilde C,\widehat D)\cap\Pi(\widehat C,\widetilde D)\\
 &\quad\text{(intersection des sous topos de }\pi(\widehat C,\widehat D)\simeq\widehat{C\times D}\text{)}
\end{aligned}
\]

% L5–9
et il est clair que les deux facteurs $\pi(\widetilde C,\widehat D)$ et $\pi(\widehat C,\widetilde D)$ sont les images inverses des sous topos $\widetilde C\subset\widehat C$ resp. $\widetilde D\subset\widehat D$ par les ``projections'' $\Pi(\widehat C,\widehat D)\longrightarrow\widehat C$ resp. $\pi(\widehat C,\widehat D)\longrightarrow\widehat D$. Cela nous ramène \unsure{à|au} \deleted{cas où} prouver que $\pi(\unsure{C|\widehat C},\unsure{\overline D|\widehat D})$ est un 2-produit des topos $\widehat C,\widehat D$. Or pour un topos \deleted{E} $X$, on a

% L10–11
\[
\begin{aligned}
\underline{\mathrm{Hom}}_{\mathrm{top}}(X,\widehat C)&\cong\underline{\mathrm{Hom}}'(C,X)^o\\
\underline{\mathrm{Hom}}_{\mathrm{top}}(X,\widehat D)&\cong\underline{\mathrm{Hom}}'(D,X)^o
\end{aligned}
\]

% L12–13
où le $\underline{\mathrm{Hom}}'$ désigne les foncteurs exacts à g [NB on a supposé $C$ et $D$ stables par $\varprojlim$ finies], donc on a

% L14–16
\[
\begin{aligned}
\underline{\mathrm{Hom}}_{\mathrm{top}}(X,\widehat C)\times\underline{\mathrm{Hom}}_{\mathrm{top}}(X,\widehat D)
 &\cong[\underline{\mathrm{Hom}}'(C,X)\times\underline{\mathrm{Hom}}'(D,X)]^o\\
 &\cong(\underline{\mathrm{Hom}}'(C\times D,X))^o\\
 &\cong\underline{\mathrm{Hom}}_{\mathrm{top}}(X,\unsure{\widehat{C\times D}|\widehat C\times D}))\qquad\text{cqfd}
\end{aligned}
\]

% L17–21
\inserted{[si $C$ et $D$ sont stables par $\varprojlim$ finies, \deleted{il} \deleted{\unsure{C}} en est de $=$ de $C\times D$, et les foncteurs exacts à gauche \inserted{(Attention, en les couples de variables, non pas en les variables séparément)} $C\times D\xrightarrow{w}X$ [$X$ une catégorie quel.] correspondent aux \deleted{\unsure{fon|deux}} couples de foncteurs exacts à gauche $C\xrightarrow{u}X$, $D\xrightarrow{v}X$, en leur associant \deleted{\unsure{$u(a)$}} $(a,b)\longrightarrow u(a)\times v(b)$, et à $w$ les foncteurs $a\longrightarrow w(a,e_D)$, $b\longrightarrow w(e_C,b\unsure{)}$}

% foot
% archive sheet number: 6
\scanpage{8}
% L1-left
\[
\text{a)}\qquad \widehat C\times\widehat C'\simeq\widehat{C\times C'}
\]

% L2-left
$C\ C'$ petites catégories avec $\varprojlim$ finies

% L3–9-left
plus généralement, si on a une famille de catégories $C_i$ petites avec $\varprojlim$ finies, soit $\prod'_i C_i\subset\prod C_i$ la sous-catégorie \inserted{pleine} du produit formé\unsure{e|{}} des systèmes $(x_i)_{i\in I}$ $(x_i\in\mathrm{Ob}\,C_i)$ tels que $x_i\simeq e_{C_i}$ pour presque tt $i$, alors

% L10-left
\[
\prod\widehat C_i\cong\widehat{\prod' C_i}
\]

% L1–2-right
\diagram{\begin{tikzcd}[ampersand replacement=\&,row sep=tiny]
C\arrow[r,"\varphi"] \& X\\
C'\arrow[ur,"\varphi'"'] \&
\end{tikzcd}}

% L1–3-right
$\varphi,\varphi'$ exacts \unsure{a|à} gauche \margin{(catégorie avec $\varprojlim$ finies)} se factorisent en

% L3–5-right
\diagram{\begin{tikzcd}[ampersand replacement=\&,row sep=small,column sep=large]
C\arrow[dr,"{i_1(x)=(x,e_{C'})}"] \& \&\\
\& C\times C'\arrow[r,"f"] \& \unsure{\mathscr X|X}\\
C'\arrow[ur,"{i_2}"',"{i_2(x')=(e_C,x')}" description,near start] \& \&
\end{tikzcd}}

% L5–6-right
par $f(x,x')\simeq\varphi(x)\times\varphi'(x')$

% L7–10-right
\[
\begin{aligned}
f(\varprojlim(x_i,x'_i))&=f(\varprojlim x_i,\varprojlim x'_i)\\
 &=\varphi(\varprojlim x_i)\times\varphi'(\varprojlim x'_i)\\
 &=\varprojlim_i\varphi(x_i)\times\varprojlim_{\unsure{i'|i}}\varphi(x'_i)\\
 &=\varprojlim_i\bigl(\varphi(x_i)\times\varphi(x'_i)\bigr)\quad\text{OK}
\end{aligned}
\]

% L11–12
b) Soient $E_i$ des topos $(i\in I)$, $E'_i\underset{\alpha_i}{\hookrightarrow}E_i$ sous-topos, qu'on \deleted{\unsure{peut}} suppose strict

% L13
Supposons $E=\prod_{i\,\mathrm{top}}E_{\unsure{j|i}}\xrightarrow{\mathrm{pr}_i}E_{\unsure{j|i}}$ existe

% L14
Soit pour tt $i$, $F_i\subset E$ le sous-topos $\mathrm{pr}_i^{-1}(E'_i)$

% L15
Soit $E'=\bigcap F_i$ leur sous-topos intersection

% L16
Donc on a \qquad $E'\xrightarrow{\unsure{\mathrm{pr}_i|\mathrm{pr}'_i}}E'_i$\qquad $E'\simeq\prod_{\mathrm{top}}E'_i$

% L17
Ceci fait de $E'$ un 2-produit des $E'_i$

% L18–19
\diagram{\begin{tikzcd}[ampersand replacement=\&,column sep=large]
\prod_i\underline{\mathrm{Hom\,top}}\,(X,E'_i)\arrow[r,hook,"{\text{pl.fid}}", "\alpha"'] \& \prod_i\underline{\mathrm{Hom\,top}}\,(X,E_i)\\
\& \underline{\mathrm{Hom\,top}}\,(X,E)\arrow[u,"{\approx\,\unsure{\beta|B}}"]
\end{tikzcd}}

% foot
% archive sheet number: 7
\scanpage{9}
% L1–7
L'image \inserted{essentielle} de $\alpha$ est la sous-catégorie \inserted{str} pleine de $\prod_i\mathrm{Hom\,top}(X,E_i)$ formée des systèmes $(f_i\colon X\longrightarrow E_i)$ tels que pour tt $i$, $f_i$ se factorise par $E'_i$ ; \unsure{ou|or} \unsure{encore|alors}, \deleted{\unsure{donc $f$}} par $\beta$, elle correspond à la ss-cat. str. pleine de $\underline{\mathrm{Hom\,top}}\,(X,E)$ formée des $f\colon X\longrightarrow E$ tels que \inserted{$\forall i$,} $\mathrm{pr}_i\circ f$ se factorise par $E'_i$, \uline{i.e.} tels que \inserted{$\forall i$,} $f$ se factorise par $F_i$, i.e. tels que $X$ se factorise par $F$. Donc

% L8
\[
\prod_i\underline{\mathrm{Hom\,top}}\,(X,E'_i)\cong\underline{\mathrm{Hom\,top}}\,(X,F).
\]

% L9–10
\uline{Corollaire de} a) et b) : Le 2-produit d'une famille de topos existe tjrs.

% L11–12
\deleted{Pour avoir la $\varprojlim_{\mathrm{top}}$ d'un topos fibré, il faut}

% L12–15-left
\diagram{\begin{tikzcd}[ampersand replacement=\&,row sep=small]
\& \deleted{\mathscr F_j} \& \deleted{\mathscr F_i}\arrow[l]\\
\deleted{\mathscr F}\arrow[d,dash] \& \deleted{\mathscr F_j}\arrow[r,"{\deleted{f^*}}"] \& \deleted{\mathscr F_i}\\
\deleted{I} \& \deleted{j} \& \deleted{i}\arrow[l,"{\deleted{f}}"]
\end{tikzcd}}

% L13–16-centre
\[
\begin{aligned}
&\deleted{\varprojlim_i\underline{\mathrm{Hom\,top}}\,(X,\mathscr F_i)}\\[-2pt]
&\deleted{\rotatebox{90}{\ensuremath{\simeq}}}\\[-2pt]
&\deleted{\varprojlim_i\underline{\mathrm{Hom}}^{\unsure{!|'}}(\mathscr F_i,X)}\\[-2pt]
&\deleted{\rotatebox{90}{\ensuremath{\simeq}}}
\end{aligned}
\]

% L15–17-right
\diagram{\begin{tikzcd}[ampersand replacement=\&,row sep=tiny,column sep=small]
\deleted{x}\arrow[r,bend left,dash] \& \deleted{y}\arrow[r,bend left,dash] \& \deleted{z}\arrow[r,"\deleted{\sim}"] \& {}\\
\& \deleted{\unsure{f^*(x)}}\arrow[u,phantom,"\simeq"] \& \deleted{g^*f^*(x)}\arrow[r,"\deleted{\sim}"']\arrow[u,phantom,"\unsure{1}"] \& \deleted{g^*(y)}\\
{}\arrow[r,"{\deleted{f}}"'] \& {}\arrow[l] \& {}\arrow[l,"{\deleted{g}}"'] \& {}\\
{}\arrow[r,dash] \& {} \& {} \& {}
\end{tikzcd}}

% L18
\deleted{$(E_i)_{i\in I}$ pseudo syst projectif de topos}

% L18–21
\diagram{\begin{tikzcd}[ampersand replacement=\&,row sep=small,column sep=small]
\& \deleted{\mathscr F_1}\arrow[r,dash,shift left]\arrow[r,dash,shift right] \& \deleted{\mathscr F_2}\arrow[r,dash,shift left=2]\arrow[r,dash]\arrow[r,dash,shift right=2] \& \deleted{\mathscr F_3}\\
\deleted{\underset{i}{\underline{\lim}}_{\mathrm{top}} E_i}\arrow[r] \& \deleted{\prod_{i\in\mathrm{Ob}\, I}E_i}\arrow[r,shift left]\arrow[r,shift right]\arrow[d,"{\deleted{\mathrm{pr}_{s(\alpha)}}}"']\arrow[dr,"{\deleted{\mathrm{pr}_{\unsure{s(\beta)|b(\alpha)}}}}"] \& \deleted{\prod_{\alpha\in\mathrm{Fl}_1 I}E_{b(\alpha)}}\arrow[r,shift left=2]\arrow[r]\arrow[r,shift right=2] \& \deleted{\prod_{(\alpha,\beta)\in\mathrm{Fl}_2(I}E_{b(\beta)}}\\
\& \deleted{E_{s(\alpha)}}\arrow[r] \& \deleted{E_{b(\alpha)}}\arrow[u,phantom,"{\deleted{\text{(but}}}"] \&
\end{tikzcd}}

% foot
% archive sheet number: 8
\scanpage{10}
% L1–4
$X,Y$ deux espaces topologiques, considérons $\mathrm{Top}(X)\mathbin{\underset{\mathrm{top}}{\overset{2}{\times}}}\mathrm{Top}(Y)$ et $\mathrm{Top}(X\times Y)$ (on vérifie facilement que ce dernier ne dépend que de $\mathrm{Top}(X)$ et de $\mathrm{Top}(Y)$, i.e. des \deleted{\unsure{les}} espaces sobres associés à $X$ et $Y$) On a un morphisme canonique de topos

% L5
\[
\deleted{\unsure{\mathrm{Top}|\mathrm{Prop.}}}\quad\tag*{\ensuremath{(*)}}\boxed{\mathrm{Top}(X\times Y)\longrightarrow\mathrm{Top}(X)\mathbin{\underset{\mathrm{top}}{\overset{2}{\times}}}\mathrm{Top}(Y)}
\]

% L6
et on veut des conditions pour que ce soit une équivalence.

% L6-left
\margin{$O_X=$ cat. des ouverts de $X$}

% L7–14
Or on voit que $\mathrm{Top}(X\times Y)=\widetilde{O_X\times O_Y}$, où la topologie à mettre sur $O_X\times O_Y$ est la ``topologie produit'' \inserted{$\pi$}, i.e. celle pour laquelle les faisceaux sont les foncteurs contravariants $(O_X\mathbin{\overset{\deleted{\unsure{o|2}}}{\times}}O_Y)^o\longrightarrow(\mathrm{Ens})$ qui sont des faisceaux \uline{séparément}. \deleted{Il} Il y a aussi la topologie $\pi'$ \deleted{\unsure{par}} ``induite'' par le produit $X\times Y$ en utilisant l'injection (supposant $X$ \deleted{\unsure{et}} $Y$ \deleted{\unsure{=}} $\ne\varnothing$) $O_X\times O_Y\subset O_{X\times Y}$. Il est clair que tout faisceau pour cette topologie (faisceau sur $X\times Y$) est un bifaisceau \deleted{\unsure{pour donc}} donc la top. $\pi'$ est plus fine que $\pi$.

% L12-left
\margin{\uline{NB} les $U\times V$ ($U\in O_X$, $V\in O_Y$) font une \uline{base} de la top. de $X\times Y$}

% L14–19
On voit déjà que le \deleted{\unsure{topos}} morphisme $(*)$ est un \uline{plongement de topos}, et on est ramené à voir si ce plongement est une équivalence i.e. $\pi=\pi'$. Donc on est amené à voir si tout recouvrement pour $\pi'$ de $U\times V$ ($U\in O_X$, $V\in O_Y$) par des $U_i\times V_i$ [i.e. rec. ensembliste dans $X\times Y$) est un rec. pour $\pi$, i.e. \deleted{si} si

% L15–end-left
% slanted
\margin{\uline{NB} Comme $\mathrm{Top}(X\times Y)$ est un sous-topos de $\mathrm{Top}(X)\times\mathrm{Top}(Y)$, ayant \uline{suffisamm}\unsure{t} de pts, et contenant tous les pts de $\mathrm{Top}(X\times Y)$, on en conclut que c'est le plus \unsure{petit} sous topos de $\mathrm{Top}(X)\times\mathrm{Top}(Y)$ contenant les pts du dit topos. Analogue pour produit d'une famille quelconque. Donc la question si \unsure{$(*)$|c'}est une équivalence \uline{équivaut} à celle si $\mathrm{Top}(X)\mathbin{\underset{\mathrm{top}}{\overset{2}{\times}}}\mathrm{Top}(Y)$ a assez de points\ldots}

% L20
\[
\bigcup_i U_i\times V_i=U\times V\quad\Longrightarrow\quad\mathop{\mathrm{Sup}}_i U_i\times V_i=U\times V
\]

% L21–22
le sup étant pris dans l'ens des sous objets de $U\times V$ dans le topos $\widetilde{O_X\times O_Y}$ (top $\pi$).

% L23–25
\uline{Proposition} \quad Pour qu'il en soit ainsi, il suffit que tout point de $U$ (\uline{ou} de $V$) admette un voisinage (pas néc. ouvert) quasi-compact

% L26–27
\uline{Corollaire} \quad \deleted{\unsure{Si}} \deleted{\unsure{tout|tous}} pt de $X$ \inserted{(ou de $Y$)} admet un syst. fond. de voisinages qui sont quasi-compacts, alors $(*)$ est une équivalence.

% foot
% archive sheet number: 9
\scanpage{11}
% L1–4
\uline{Dém} \quad Pour tt $x\in X$, soit $U_x$ un voisinage qu. cpt. Donc $\forall y\in Y$ il existe $I_{x,y}\subset I$, $I_{x,y}$ partie finie, telle que \deleted{\unsure{$y$}} $U_x\times\{y\}$ soit recouvert par les $U_i\times V_i$ ($i\in I_{xy}$). Soit $V_{x,y}$ l'intersection \deleted{de \unsure{ces|des} $V_i$} $V_{x,y}=\bigcap_{i\in I_{xy}}V_i$, donc on a

% L5–6
\[
\left\{\begin{aligned}
\mathring U_x=U'_x&\subset\bigcup_{i\in I_{xy}}U_i,\\
V_i&\supset V_{x,y}\quad\text{pour }i\in I_{xy}
\end{aligned}\right.
\qquad\deleted{\unsure{\mathring U_x\times V_{x,y}\subset\bigcup_{i\in I_{xy}}U_i\times V_{xy}}}
\]

% L7
\[
\begin{aligned}
\text{Donc on a}\quad S\overset{\mathrm{déf}}{=}\mathop{\mathrm{Sup}}_{i\in I}(U_i\times V_i)
 &\geq\mathop{\mathrm{Sup}}_{i\in I_{x,y}}\deleted{\unsure{U_x\times V_{xy}}}(U_i\times V_i)\\
 &\geq\inserted{\mathop{\mathrm{Sup}}_{i\in I_{xy}}(U_i\times V_{xy})=}
 \bigl(\mathop{\mathrm{Sup}}_{i\in I_{xy}}U_i\bigr)\times V_{xy}\geq U'_x\times V_{xy}
\end{aligned}
\]

% L8
donc comme ceci est vrai pour tt $y$, on a

% L9–10
\[
\begin{aligned}
 S&\geq\mathop{\mathrm{Sup}}_y\mathring U'_x\times V_{xy}
 =\mathring U'_x\times\mathop{\mathrm{Sup}}_y V_{xy}=\mathring U'_x\times V\\
 \text{donc}\quad S&\geq\mathop{\mathrm{Sup}}_x(\mathring U'_x\times V)
 =\bigl(\mathop{\mathrm{Sup}}_x\mathring U'_x\bigr)\times V=U\times V
\end{aligned}
\]

% L11
\hfill cqfd.

% L12–14
\uline{Raffinement} \quad Supposons que $X$ \inserted{(ou $Y$)} ne contienne pas de partie \inserted{loc. fermée} non vide dont aucun point n'ait de voisinage \inserted{relatif} quasi-compact. Je dis qu'alors $(*)$ est une équivalence

% L15–21
Considérons en effet les ouverts $U'$ dans $U$ tels que $S\geq U'\times V$, il suffit de prouver que $U$ est parmi. Or la \deleted{\unsure{partie des ouverts}} \inserted{réunion de ces} ouverts est du $=$ type, donc il y a un plus grand ouvert \inserted{$U'$} de ce type [N.B $\unsure{\varnothing_X|\varnothing x}$ en est !] prouvons $U'=U$. \deleted{Soit $Z$} Sinon, la partie loc. fermée $Z$ de $U$ \deleted{\unsure{serait}} serait non vide, donc on aurait un $x\in Z$ ayant un voisinage relatif $T_{\unsure{x|*}}\subset Z$ quasi-compact. Soit $U''\subset U$ un ouvert contenant $x$ tel que $U''\cap Z\subset T$. \deleted{\unsure{On a|Donc}} Il

% L22–23
existe $\forall y\in Y$ une partie finie $I_y$ de $I$ telle que $T\times\{y\}\subset\bigcup_{i\in I_y}U_i\times V_i$, et posant $V_y=\bigcap_{i\in I_y}V_i$, \deleted{$U_x=\bigcap_{i\in\unsure{I_y}}U_i$}, \deleted{\unsure{on aura}}, on

% L24
\[
\deleted{T\times V_y\subset\bigcup_{i\in\unsure{I}}U_i\times V_i\quad\text{donc}}\quad
 T\subset\bigcup_{i\in I_{\unsure{y|xy}}}U_i,\qquad i\in I_y\Rightarrow V_i\supset V_y
\]

% L25
donc on a

% L26
\[
S\geq\mathop{\mathrm{Sup}}_{i\in I_y}(U_i\times V_i)\geq\mathop{\mathrm{Sup}}_{i\in I_y}U_i\times V_y\geq\bigl(\mathop{\mathrm{Sup}}_{i\in I_y}U_i\bigr)\times V_y
\]

% L27–28
\[
\begin{aligned}
\text{Mais on sait déjà que}\qquad &S\geq U'\times V\geq U'\times V_y\\
\text{donc}\qquad &S\geq\mathrm{Sup}(U',U_i\ i\in I_y)\times V_y\geq(U'\vee U'')\times V_y
\end{aligned}
\]

% foot
% archive sheet number: 10
\scanpage{12}
% L1
Comme ceci est vrai pour tt $y$, on trouve

% L2
\[
S\geq\mathrm{Sup}\bigl((U'\cup U'')\times V_y\bigr)
 =(U'\cup U'')\times\mathrm{Sup}\,V_y=(U'\cup U'')\times V
\]

% L3
contrairement au caractère maximal de $U'$, OK.

% L4
\uline{Cas de produits \deleted{\unsure{infinis|infin.}} d'une famille quelconque de topos}

% L5–8
On trouve que si tous les facteurs sauf \inserted{au plus} un seul sont tels que tt pt ait un système fond. de voisinages (pas néc. ouverts) quasi-compacts, et si tous sauf au plus un seul, sont \deleted{\unsure{munis}} eux-mêmes quasi-compacts, alors

% L9
\[
\mathrm{Top}\bigl(\prod_{i\in I}X_i\bigr)\longrightarrow\prod_{i\,(\mathrm{top})}^{\unsure{(2)|(b)}}\mathrm{Top}(X_i)
\]

% L10
est une équivalence de topos.

% L11–14
\deleted{C'est essentiellement le $=$ argument, travaillant sur la cat. \unsure{du} produit restreint $\prod' O_{X_i}$ (familles \deleted{\unsure{de|des}} $(U_i)_{i\in I}$ $U_i\in O_{X_i}$ pour tt $i$, $U_i=X_i$ (\inserted{plus} \deleted{\unsure{objet}} grand objet de $O_{X_i}$) pour presque tout $i$) et la ``topologie produit'' sur ce produit.}

% L15–17
Dans le cas où $I$ est fini, c'est un corollaire immédiat de ce qui précède. Dans le cas général, on a $I=\varinjlim_\alpha I_\alpha$ ($I_\alpha$ parties finies de $I$) donc

% L18–19
\[
\begin{aligned}
\prod_{i\,\mathrm{top}}\mathrm{Top}(X_i)
 &\cong\mathop{\varprojlim_{\mathrm{top}}}_{\alpha}\bigl(\prod_{i\in I_\alpha}\mathrm{Top}(X_i)\bigr)
 \cong\mathop{\varprojlim_{\mathrm{top}}}_{\alpha}\bigl(\mathrm{Top}(\underbrace{\prod_{i\in I_\alpha}X_i}_{Z_\alpha}\unsure{)}\bigr)\\
 &\cong\bigl(\underbrace{\varinjlim_\alpha O_{Z_\alpha}}_{O'_Z}\bigr)^{\sim}
\end{aligned}
\]

% L20–26
$O'_Z$ est la catégorie des ouverts de $Z$ \inserted{$=\prod X_i$} qui sont images inverses d'ouverts de produits finis $\prod_{i\in I_\alpha}X_i=Z_\alpha$. \deleted{\unsure{Ils}} On \deleted{\unsure{peut|une}} a dessus la topologie \inserted{$\pi$ des sites} $\varinjlim$, qui \deleted{\unsure{définit}} est utilisée dans la formule précédente, mais aussi la top. plus fine $\pi'$ provenant de la top de $Z$. Si tous les \deleted{\unsure{Z}} \unsure{$X_i$|$i$} sauf un nb fini sont quasi-compacts [les autres étant ce que ils veulent] on trouve par un argument tt analogue \inserted{à dessus} que les top. sont

% foot
% archive sheet number: 11
\scanpage{13}
% L1–3
les \unsure{$\widetilde{\mathrm m}$s|mêmes}. \deleted{Soit en effet} \inserted{$W\in O'$, donc de la forme $U\times Z'_\alpha$,} \deleted{$\unsure{\prod_{i\in I_\alpha}U_i}\times\prod_{j\in I-I_\alpha}X_j$} $U$ un ouvert de $Z_\alpha$, et $Z'_\alpha=\prod_{i\in I-I_\alpha}X_i$, (d'où $Z=Z_\alpha\times Z'_\alpha$), et soient

% L4–6
$(V_\lambda)_{\lambda\in\Lambda}$, $V_\lambda\in O'_Z$, avec $W\subset\bigcup_\lambda V_\lambda$. prouvons qu'\unsure{on|ils|il} \deleted{\unsure{existe}} \deleted{\unsure{t}} a $\mathop{\mathrm{Sup}}_\lambda W\geq\inserted{W=}U\times Z'_\alpha$. \unsure{OBS|OPS} \deleted{\unsure{On}} $I_\alpha$ assez grand pour que les $X_i$ ($i\in I-I_\alpha$) soient quasi-cpts, donc $Z'_\alpha$ qu. cpt.

% L7–8
$\forall z\in\deleted{\unsure{Z_\alpha}} W$, $\{z\}\times Z'_\alpha$ est recouvert par un nb fini de $V_\lambda$,

% L8–12
\deleted{disons pour $\lambda\in\Lambda_z$, donc $\exists\unsure{I_{\beta(z)}}$ \unsure{à} $I_{\beta(z)}\supset I_\alpha$ \unsure{--} tels que chaque $V_\lambda$ ($\lambda\in\Lambda_z$) de la forme $V'_\lambda\times Z'_{\beta(z)}$, avec $V'_\lambda\in O_{Z_{\beta(z)}}$. \unsure{$V_\lambda$} Comme $\underset{\{z\}\times Z_{\alpha,\beta(z)}\times Z'_{\beta(z)}}{\{z\}\times Z'_\alpha}\subset\bigcup_\lambda(V'_\lambda\times Z'_{\beta(z)})$, cela \deleted{\unsure{est faux que} $Z'$} signifie $\{z\}\times\prod_{i\in I_{\beta(z)}-I_\alpha} X_i\subset\bigcup_{\lambda\in\Lambda_z}V'_\lambda$, donc il existe}

% L13–15
donc \deleted{\unsure{on a}} $\exists$ un voisinage ouvert $U_z\ni z$ dans $Z_\alpha$ tel que $U_z\times Z'_\alpha\subset\bigcup_{\lambda\in\Lambda_z}V_\lambda$. Mais alors \deleted{\unsure{c'est une}} on aura l'inclusion $U_z\times\prod_{i\in I_{\beta(z)}-I_\alpha}X_i\subset\bigcup_{\lambda\in\Lambda_z}V'_\lambda$ donc

% L16
\[
S\overset{\mathrm{déf}}{=}\mathop{\mathrm{Sup}}_{\lambda\in\Lambda\deleted{\unsure{z}}}(V_\lambda)\geq\mathop{\mathrm{Sup}}_{\lambda\in\Lambda_z}V_\lambda\geq U_z\times Z'_\alpha.
\]

% L16–17
Comme ceci est vrai pour tt $z$, on a \deleted{\unsure{Sup}}

% L17–18
\[
\begin{aligned}
 S&\geq\mathop{\mathrm{Sup}}_z\bigl(U_z\times Z'_\alpha\bigr]=\bigl(\mathop{\mathrm{Sup}}_{z\in\unsure{Z|z}}U_z\bigr)\times Z'_\alpha\\
 &=W\times Z'_\alpha,\quad\text{cqfd.}
\end{aligned}
\]

% L19
\figure{\begin{tikzpicture}\draw (0,0)--(2,0);\draw (.35,-.13)--(1.75,-.13);\end{tikzpicture}}

% L20
Contr\unsure{'|-}exemples pour des produits \uline{infinis}

% L21
Prenons $X_i$ discret pour tout $i$, on aura

% L22
\[
\prod_{i\in I}\mathrm{Top}(X_i)=\underbrace{\widetilde{\varinjlim_\alpha O_{Z_\alpha}}}_{O'_Z},\quad\text{montrons que la top.}
\]

% foot
% archive sheet number: 12
\scanpage{15}
% L1–3
$\pi$ de $O'_Z$ n'est \uline{pas} la topologie déduite de la top produit sur $Z=\unsure{\prod\deleted{X_i}{}^I_i|\prod X_i|\deleted{\prod X_i}}=\varprojlim Z_\alpha$ (si $I$ infini et les $X_i$ ne sont pas presque tous finis).

% L3–7
\deleted{\unsure{En d'autres termes|lim d'ouverts fermés} Montrons en effet qu'on peut trouver une fam} On a une famille fond. de \unsure{recouvrements|voisinages} \deleted{\unsure{en}} de $Z$ par des objets de $O'$ en associant à tt $z=(x_i)_{i\in I}\in Z$ une partie finie $J(z)$ de $I$, \deleted{telle} \inserted{posant} \deleted{que \unsure{$\exists W(z)=$}} \inserted{alors} $V_z=\{z'\in Z\mid z'_i=z_i\ \text{si }i\in I-J(z)\}$.

% L8–9-left
\margin{?}

% L8–9
\deleted{La question est de voir si tous ces} $(V_z)_{z\in Z}$ est \inserted{tjrs} un rec. pour $\pi$, ou encore

% foot
% archive sheet number: 14
\scanpage{16}
% L1
\uline{$\varprojlim$ filtrantes de topos essentiels} $\widehat C$

% L2–4
On est intéressé aux $\varprojlim$ pour des topos fibrés où les fibres sont des topos $\widehat C$ ($C$ karoubienne \inserted{petite} si on veut) et les morphismes de transition \inserted{$f\colon\widehat C'\longrightarrow\widehat C$} sont essentiels, i.e. associés à des foncteurs

% L5
\[
f_{\unsure{0|o}}\colon C\longrightarrow C'.
\]

% L6–7
Il est douteux qu'alors \inserted{$\unsure{\mathscr X|\mathcal X}=$} $\varprojlim_{\mathrm{top}}\widehat C_i$ soit un topos essentiel et ses projections $\mathscr X\longrightarrow\widehat C_i$ des morphismes essentiels.

% L8–10
Mais pour le calcul commode de $\mathscr X$ on va supposer que les $\widehat f_{ij}\colon\widehat C_j\longrightarrow\widehat C_i$ \deleted{\unsure{soient}} (associés au $f_{ij}\colon C_j\longrightarrow C_i$) soient ass. à des morph. de sites donnés par des \unsure{foncteurs|$\infty$-foncteurs}

% L11
\[
g_{\unsure{ji|ij}}\colon C_i\longrightarrow C_j
\]

% L12
(préc. induits par $\widehat f_{ij}^{*}$ à iso. près, et adjoints a \uline{droite} des $f_{ij}$).

% L13
On trouve alors

% L14
\[
\mathscr X:\mathop{\varprojlim_{\mathrm{top}}}_{i,\widehat f_{ij}}\widehat C_i
\simeq\bigl(C=\mathop{\underline{\lim}}_{i,g_{ij}}C_i\bigr)^{\wedge}
\]

% L15–18
et les morphismes $\mathscr X\xrightarrow{\mathrm{pr}_i}\widehat C_i$ sont associés aux morphismes de sites donnés par les foncteurs can $C_i\xrightarrow{\alpha_i}C=\mathop{\varinjlim}_j C_j$, donc pour $F\in\widehat C=\mathscr X$, on a $\mathrm{pr}_{i*}(F)=F\circ\alpha_i$, et pour $G\in\widehat C_i$ on a

% L19
\[
\mathrm{pr}_i^*(\unsure{\widehat G|G})=
\]

% foot
% archive sheet number: 15
\scanpage{18}
% L1
foncteurs sur un topos $\widetilde C$ commutant aux $\varinjlim$

% L2–3
Soient $C$ un $\mathcal U$-site, $E$ une catégorie avec petites $\varinjlim$, considérons les catégories

% L4
\[
\underline{\mathrm{Hom}}'(\widetilde C,E)\qquad\underline{\mathrm{Hom}}\,(C,E)
\]

% L5–6
\deleted{et} \deleted{Soit} $\underline{\mathrm{Hom}}'(\widetilde C,E)$ \deleted{désignant} la sous-catégorie pleine de $\underline{\mathrm{Hom}}$ formée des foncteurs qui commutent aux petites $\varinjlim$. On va définir des foncteurs

% L7
\diagram{\begin{tikzcd}[ampersand replacement=\&,column sep=large]
\underline{\mathrm{Hom}}\,(\widetilde C,E)\arrow[r,shift left,"\alpha"] \& \underline{\mathrm{Hom}}\,(C,E)\arrow[l,shift left,"\beta"]
\end{tikzcd}}

% L8
en posant

% L9–10
\[
\begin{aligned}
\alpha(u)&=u\circ\varepsilon_C
 &&(\varepsilon_C\colon C\to\widetilde C\ \text{le foncteur canonique})\\
\beta(u_0)&=[F\to\mathop{\underline{\lim}}_{C/F}u_0(x)]
 &&(=\widehat\beta\mid\widetilde C,\ \text{où }\widehat\beta\colon\widehat C\longrightarrow E\ \text{est}\\
&&&\text{le prolongement can. de }\beta\ \text{défini}\\
&&&\text{par commutation aux }\varinjlim)
\end{aligned}
\]

% L11–14-right
\diagram{\begin{tikzcd}[ampersand replacement=\&,row sep=small,execute at end picture={\draw (current bounding box.north east) -- (current bounding box.north west) -- (current bounding box.south west) -- (current bounding box.south east);}]
C\arrow[r,hook,"{j_0}"'] \& \widehat C\arrow[dl,bend left,"a"] \& \varepsilon_C=aj_0\\
\widetilde C\arrow[ur,hook,"{i_0}"] \& \&
\end{tikzcd}}

% L11–12
Soit $u\in\mathrm{Ob}\,\underline{\mathrm{Hom}}\,(\unsure{\overline C|\widetilde C},E)$, on a un hom

% L13
\[
i_u\colon\beta\alpha(u)\longrightarrow u
\]

% L14
donné \unsure{pour|par} $\forall F\in\mathrm{Ob}\,\widetilde C$ par

% L15–16
\diagram{\begin{tikzcd}[ampersand replacement=\&,row sep=small]
{i_u(F)\colon} \& (\beta(\alpha(u)))(F)\arrow[r] \& u(F)\\
\& \varinjlim_{C/F}u(\varepsilon_C(x))\arrow[u,equal] \&
\end{tikzcd}}

% L16–18-right
déduit de l'isom can. $\mathop{\varinjlim}_{C/F\,(\widetilde C)}X\xrightarrow{\sim} F$ en appliquant $u$

% L19–20
On trouve que l'hom $i_u(F)$ est bien fonctoriel en $F$, d'où $i_u$, et que $i_u$ est fonctoriel en $u$, d'où

% L21
\[
\boxed{\beta\alpha\xrightarrow{i}\mathrm{id}_{\underline{\mathrm{Hom}}(\widetilde C,E)}}
\]

% L22–24
Notons tout de suite que \deleted{\unsure{$u$}} $i_u\colon\beta\alpha(u)\longrightarrow u$ est un isom \unsure{sss|ssi} $u\colon\widetilde C\longrightarrow E$ commute aux $\varinjlim$ du type $\varinjlim_{C/F}X\xrightarrow{\sim}F$ ($F\in\mathrm{Ob}\,\widetilde C$) ; Il \uline{suffit} \inserted{donc} que $u\in\mathrm{Ob}\,\underline{\mathrm{Hom}}'(\widetilde C,E)$.

% L25
Soit $u_0\in\underline{\mathrm{Hom}}\,(C,E)$, on va définir

% L26
\[
j_{u_0}\colon\deleted{\unsure{\alpha\beta}}u_0\deleted{\unsure{\}}}\unsure{\longleftrightarrow|\longrightarrow}\alpha\beta(u_0)
\]

% L27
\deleted{en posant}, \inserted{donné} pour $Z\in\mathrm{Ob}\,C$, par

% L28–30
\diagram{\begin{tikzcd}[ampersand replacement=\&,row sep=small]
{j_{\unsure{u_0|0}}(Z)\colon}\ \deleted{\unsure{(\alpha(\beta(u_0)))(Z)\longrightarrow}}
 \& u_0(Z)\arrow[r] \& (\alpha(\beta(u_0)))(Z)\\
\& \& \beta(u_0)(\varepsilon_C(Z))\arrow[u,equal]\\
\& \& \varinjlim_{C/\varepsilon_C(Z)}u_0(X)\arrow[u,equal]
\end{tikzcd}}

% L31–32
en notant que $(Z,\deleted{\unsure{\varepsilon_C(Z)}}\inserted{\text{morphisme can. de } Z\text{ dans }\varepsilon_C(Z)}\unsure{\}})$ est un objet de $C/\varepsilon_C(Z)$, soit $X_0$, et que par suite $u_0(X_0)=u_0(Z)$ s'en va dans la $\varinjlim$.

% foot
% archive sheet number: 17
\scanpage{19}
% L1–3
On vérifie que cette application est fonctorielle en $Z$, donc définit $j_{u_0}\colon u_0\longrightarrow\alpha\beta(u_0)$, et que cette dernière est fonctorielle en $u_0$, d'où un hom de foncteurs

% L4
\[
\boxed{j\colon\mathrm{id}_{\underline{\mathrm{Hom}}(C,E)}\longrightarrow\alpha\beta}
\]

% L5–6
Signalons tout de suite que pour \deleted{\unsure{$Z$|que}} $u_0$ donné, $j_{u_0}\colon u_0\longrightarrow\alpha\beta(u_0)$ est un isom. sss \deleted{\unsure{on}} $\forall Z\in\mathrm{Ob}\,C$,

% L7
\[
u_0(z)\to\varinjlim_{C/\varepsilon_C(z)}u_0(X)\qquad\text{est un isom,}
\]

% L8–9
ou encore sss pour tout tel $Z$ et $\xi\in\mathrm{Ob}\,E$, la flèche déduite de la précédente en appliquant $\mathrm{Hom}(-,\xi)$ est bijective,

% L10
\[
\text{i.e.}\qquad\varprojlim_{C/\varepsilon_C(z)}\mathrm{Hom}_{\unsure{\iota|1|{}}}(u_0(X),\xi)\longrightarrow\mathrm{Hom}(u_0(z),\xi)\qquad\text{bijective.}
\]

% L11
Or soit \deleted{\unsure{pour $\xi\in\mathrm{Ob}\,E$}},

% L12
\[
u_0^{\unsure{\bullet|\cdot|\prime}}\colon E\longrightarrow\widehat C
\]

% L13
le foncteur défini par

% L14
\[
u_0^\bullet(\xi)(Y)=\mathrm{Hom}(u_0(Y),\xi)
\]

% L15
de sorte que l'application précédente devient, posant $P=u_0^\bullet(\xi)$ :

% L16
\[
\varprojlim_{C/\varepsilon_C(z)}P(X)\longrightarrow P(Z)\qquad\text{un isom.}
\]

% L17–19
\deleted{\unsure{Nous proposons}} Cette condition est vide si $\varepsilon_C$ est pl. fidèle i.e. la top. de $C$ est moins fine que la top. canonique (On sait que $C/\varepsilon_C(z)\simeq C/z$ admet $Z$ comme objet final).

% L20–24
On va prouver que $i,j$ sont des données d'adjonction associées pour les foncteurs $\alpha$ et $\beta$, [\unsure{Celà|Cela|Ce} qui prouve que $\alpha$ et $\beta$ \deleted{établissent} \inserted{induisent} des équivalences entre les sous-catégories \inserted{str. pleines} de $\underline{\mathrm{Hom}}\,(\widetilde C,E)$ et $\underline{\mathrm{Hom}}\,(C,E)$ qu'on a décrit \unsure{auparavant|au par. n$^\circ$ 7|au préal}. On \deleted{\unsure{en conclut}} \inserted{en conclut} \deleted{\unsure{a établi}} que $\alpha$ induit un foncteur \uline{pl. fidèle}

% L25
\[
\underline{\mathrm{Hom}}'(\widetilde C,E)\longrightarrow\underline{\mathrm{Hom}}\,(C,E)
\]

% L26–28
dont l'image essentielle est formée des \deleted{\unsure{$Z$}} $u_0\colon C\longrightarrow E$ satisfaisant la condition ci-dessus, et tels que $\beta(u_0)\in\underline{\mathrm{Hom}}'(\widetilde C,E)$ -- condition qu'on remplacera par une autre plus simple]

% foot
% archive sheet number: 18
\scanpage{20}
% L1–2
$C$ un $\mathcal U$-site, $E$ catégorie où les petites $\mathop{\underline{\lim}}$ sont repr. déterminer les foncteurs

% L3
\[
u\colon\widetilde C\longrightarrow E
\]

% L4–5
commutant aux $\mathop{\underline{\lim}}$ (petites) ? Connu \inserted{à isom près} quand on connait $u\circ\varepsilon_C=u_0$

% L6
\[
u_0\colon C\longrightarrow E
\]

% L6-right
\diagram{\begin{tikzcd}[ampersand replacement=\&]
C\arrow[r,"{\varepsilon_C}"]\arrow[rr,bend right=27,"{u_0}"'] \& \widetilde C\arrow[r,"u"] \& E
\end{tikzcd}}

% L7
car pour $F\in\widetilde C$, on a un isom. fonctoriel en $F$

% L8
\[
F=\varinjlim_{C/F}\varepsilon_C(x)
\]

% L9
\[
\text{d'où}\qquad u(F)=\mathop{\underline{\lim}}_{C/F}u_0(X).
\]

% L10
De façon précise, on trouve un foncteur $u\longmapsto u\circ\varepsilon_C$

% L11
\[
\alpha\colon\underline{\mathrm{Hom}}'(\widetilde C,E)\longrightarrow\underline{\mathrm{Hom}}\,(C,E)
\]

% L12–14
qui est \uline{pl. fid}, le signe $\underline{\mathrm{Hom}}'$ désignant les foncteurs commutant aux petites $\varinjlim$. En fait, on a un foncteur ``prolongement à $\widetilde C$''

% L14
\[
\beta\colon\underline{\mathrm{Hom}}\,(C,E)\longrightarrow\underline{\mathrm{Hom}}\,(C,E)
\]

% L15
défini par

% L16
\[
\beta(u_0)=F\longmapsto\mathop{\underline{\lim}}_{C/F}u_0(X)
\]

% L17
et on a un isom canonique

% L18
\[
\beta\unsure{\circ|\xi}\alpha\simeq i\quad
(\text{inclusion }\deleted{\unsure{\text{de}}}\ \underline{\mathrm{Hom}}'(\widetilde C,E)\hookrightarrow\underline{\mathrm{Hom}}\,(\unsure{\overline C|\widetilde C},E))
\]

% L19–21
Mais il n'est pas clair que $\forall u_0\colon C\longrightarrow E$, $\beta(u_0)$ soit dans $\underline{\mathrm{Hom}}'$. A vrai dire, $\beta(u_0)$ est induit par $\widehat\beta(u_0)\colon\widehat C\longrightarrow E$ défini de la $=$ façon, et on sait que $\widehat\beta(u_0)$ commute aux $\varinjlim$, en fait,

% L22
\diagram{\begin{tikzcd}[ampersand replacement=\&,column sep=large]
\underline{\mathrm{Hom}}\,(C,E)\arrow[r,shift left,"{\cong}","{\widehat\beta}" description,near start] \& \underline{\mathrm{Hom}}'(\widehat C,E)\arrow[l,shift left,"{\widehat\alpha}"]
\end{tikzcd}}

% L22–24
(cf rappel plus bas). Identifions dorénavant (modulo équivalence) $\underline{\mathrm{Hom}}\,(C,E)$ à $\underline{\mathrm{Hom}}'(\widehat C,E)$, alors $\alpha$ et $\beta$ deviennent des foncteurs $\alpha'$, $\beta'$

% L25–26
\[
\begin{aligned}
\underline{\mathrm{Hom}}'(\widetilde C,E)&\xrightarrow{\alpha'}\underline{\mathrm{Hom}}'(\widehat C,E)\\
\underline{\mathrm{Hom}}'(\widehat C,E)&\xrightarrow{\beta'}\underline{\mathrm{Hom}}\,(\widetilde C,E)
\end{aligned}
\]

% foot
% archive sheet number: 19
\scanpage{21}
% L1–3
qui s'explicitent ainsi : \deleted{\unsure{soit}} $a\colon\widehat C\longrightarrow\widetilde C$ le foncteur ``faisceau associé'', donc $\varepsilon_C=a\,i_C$ ($i_C\colon C\longrightarrow\widehat C$) donc pour $u\in\underline{\mathrm{Hom}}(\unsure{\overline C|\widetilde C},E)$, $u\circ\varepsilon_C=(ua)i_C$, donc (\deleted{\unsure{compte tenu}} que $a$, donc $ua$, commute aux petites $\varinjlim$)

% L4
\[
\alpha'(\unsure{u})=u\circ a
\]

% L5
\deleted{Donc} D'autre part, on voit que pour $v\in\underline{\mathrm{Hom}}'(\widehat C,E)$ (correspondant à $v_0\colon C\longrightarrow E$)

% L6
\[
\beta'(v)=v\circ j_{\unsure{C|a}}=v\mid\widetilde C\qquad(\text{où }j\colon\widetilde C\longrightarrow\widehat C\text{ est l'inclusion})
\]

% L7
\[
\text{On a bien}\qquad\beta'\alpha'(u)=u\circ a\mid\widetilde C\simeq u
\]

% L8–12
Le fait que $\alpha'$ est pl. fidèle provient du fait que $\widetilde C$ est, moyennant $a$, une catégorie de fractions de $\widehat C$, de sorte que $\underline{\mathrm{Hom}}(\widetilde C,E)$ (sans $'$) \deleted{est} s'identifie à la sous catégorie pleine de $\underline{\mathrm{Hom}}(\widehat C,E)$ formée des foncteurs \inserted{$F$} qui sont tels que \deleted{\unsure{$F(u)$}} $u\in\mathrm{Fl}\,\widehat C$, $a(u)$ inversible $\Longrightarrow F(u)$ inversible. Est-il vrai
% L13–15
que dans cette identification, les foncteurs qui commutent aux $\varinjlim$ petites se correspondent ? Il reste à voir que si $\unsure{H|U|\mathcal U}\colon\widetilde C\longrightarrow E$ est tel que $H\circ a$ commute aux $\varinjlim$, alors
% L16–19
$H$ aussi, or on aura $H=(H\circ a)i$, et si on a des objets $X_\alpha$ de $\widetilde C$, $\varinjlim_\alpha\,{}_{\widetilde C}X_\alpha=a(\varinjlim_\alpha\,{}_{\widehat C}X_\alpha)$, donc $H(\varinjlim_\alpha\,{}_{\widetilde C}X_\alpha)=H(a(\varinjlim_\alpha\,{}_{\widehat C}X_\alpha)$

% L20
\[
=(H\circ a)(\mathop{\underline{\lim}}_\alpha X_\alpha)\underset{\substack{\text{par hyp}\\\text{sur }H\circ a}}{=}\mathop{\underline{\lim}}_\alpha(H\circ a)(X_\alpha)=\varinjlim_\alpha H(X_\alpha)\qquad\text{OK.}
\]

% L21–24
Donc on trouve que $\alpha'$ induit

% L21–24-right
\[
\begin{aligned}
\underline{\mathrm{Hom}}'(\widetilde C,E)\xrightarrow{\simeq}{}&
\text{sous-catégorie \inserted{pleine} de }\underline{\mathrm{Hom}}^{\unsure{F|{}}}(\widehat C,E)\text{ formée des}\\
&\text{foncteurs qui a) commutent aux }\mathop{\underline{\lim}}\text{ petites}\\
&\hphantom{\text{foncteurs qui }}\text{b) Transforme morphismes}\\
&\hphantom{\text{foncteurs qui }}\text{``bicouvrants'' de }\widehat C\text{ (i.e. tels que}\\
&\hphantom{\text{foncteurs qui }}a(u)\text{ inv.) en isom.}
\end{aligned}
\]

% L25
Donc on trouve

% L26–29
\[
\begin{aligned}
\underline{\mathrm{Hom}}'(\widetilde C,E)\xrightarrow[\simeq]{\alpha_0}{}&
\text{sous-catégorie pleine de }\underline{\mathrm{Hom}}(C,\unsure{E|\deleted{\widetilde E}E})\\
&\text{formée des }u_0\colon C\longrightarrow E\text{ tels que}\\
&\widehat u_0\colon\widehat C\longrightarrow E\text{ transforme morphismes}\\
&\text{``bicouvrants'' en isom.}
\end{aligned}
\]

% L30–32
Mais cette condition implique que \deleted{\unsure{la|tt}} famille couvrante \inserted{$X_\alpha\longrightarrow X$} de $\unsure{\overline C|C|\widetilde C}$ (définissant donc une famille épimorphique $\varepsilon_C(X_\alpha)\longrightarrow\varepsilon_C(X)$ de $\widetilde C$) \deleted{est} transformée en famille épimorphique de $E$.

% L33
Cela \uline{est-il suffisant} ? Supposons cette condition remplie,

% L33-left
\margin{\uline{non}}

% foot
% archive sheet number: 20
